% =============================================================================
\chapter{M1b: the stop rule and the free screen tie}\label{ch:m1b}

\begin{ideia}[M1b in two sentences]
A stop rule uses the radius to decide when to widen the report; a free screen (leave a claim out,
re-close, see whether the set changes) decides almost the same thing without computing any radius.
M1b measured whether the rule beats such a screen, and at matched depth it does not.
\end{ideia}

Table~1 of Chapter~\ref{ch:gamma} compares intervals that always widen with a narrow committed
report. It never shows the measure being \emph{used}. The table design of 7 September had asked for
exactly that row, and I had not reread it when the outline was rebuilt on 12 September. M1b is that
row, run on the M1 harness in the night of 13 September, with five predictions registered in
Appendix~B of \texttt{results/interval/PREREGISTRATION.md} before any policy quantity was computed.
M1's own cells did not move: projected onto M1's columns the new \texttt{rows.csv} has the same hash.

\begin{antes}
\antesq[tabelas-pagas]{Which comparison did the 7 September design fix in advance as the only one a
radius could win against a free statistic?}
\antesq[tabelas-pagas]{What sentence was written, before measuring, for the case of a tie?}
\end{antes}

\section{What the 7 September design asked for}

Three conclusions of that design are the ones M1b tested. \textbf{Rows are complete procedures}, and
the natural use of the radius in that format is a stop rule: report the committed interval
when the set survives every retraction within the budget, widen only when it does not. \textbf{A free
screen is the same policy at one claim}: retract one claim, re-close under Meek, flag if the committed
set changes. An unchanged set under every single retraction is a set that stays valid under every
single retraction, so at depth one the screen cannot certify what the rule refuses. The only place
the radius could beat a free statistic is at two wrong claims, where two retractions interact; the
design fixed that comparison in advance (the starred comparison) and wrote the caption for the
null result before any number existed: the work would then be ``a measurement of when knowledge
error is silent rather than a method''. \textbf{The price of the guarantee is a cell of its own}: the
share of correct analyses the rule refuses when the knowledge is right.

\begin{armadilha}[Two things called ``stop rule'' in my September notes]
My notes of 12 September (in Portuguese) call the ledger's halt-and-inspect clause, the one that found
the bucket bugs and the back-door oracle of Chapter~\ref{ch:table4}, a \emph{regra de parada}. Here that
clause is the \textbf{inspection trigger}; the \textbf{stop rule} is the reporting policy
$\mathrm{STOP}_r$.
\end{armadilha}

\prompt[conceito=tabelas-pagas,dificuldade=0.7]{Which comparison did the 7 September design fix in
advance, why only that one, and what sentence did it write for a tie?}{The two-claim stop rule
against the leave-one-out screen, at two wrong claims. Only that one, because at one claim the screen
and the rule are the same policy (a set unchanged by every single retraction is valid under every
single retraction), so any difference has to come from two retractions interacting. The sentence: the
work would be ``a measurement of when knowledge error is silent rather than a method''.}

\begin{antes}
\antesq[screen-leave-k-out]{In what class of cases do the stop rule and the one-claim screen decide
differently?}
\antesq[regra-de-parada]{Why is ``refusal implies trigger'' an argument plus a measurement, and not a
theorem?}
\antesq[orcamento-retracao]{Why does comparing the two-claim rule with the one-claim screen mix two
things?}
\end{antes}

\section{Four policies, one argument, five predictions}

\begin{exemplo}[Three retractions, two decisions (schematic)]
$\Gz$ commits to $Z$ and holds three claims; what each single retraction does, and what each policy
decides:

\begin{center}
\begin{tikzpicture}[node distance=0.22cm and 0.9cm,
  box/.style={blk, font=\scriptsize, text width=4.3cm, align=left, inner sep=3pt},
  dec/.style={font=\scriptsize\sffamily, align=left, text width=4.6cm}]
\node[blkc, font=\scriptsize, align=center, inner sep=4pt] (g0) {$\Gz$ commits to $Z$\\ claims $c_1, c_2, c_3$};
\node[box, right=1.1cm of g0, yshift=0.95cm] (r1) {retract $c_1$: nothing changes};
\node[box, draw=vermesc, fill=vermesc!7, below=of r1] (r2) {retract $c_2$: $Z$ is invalid in the re-closure};
\node[box, draw=verdesc, fill=verdesc!7, below=of r2] (r3) {retract $c_3$: the optimal set becomes $Z' \neq Z$, and $Z$ is still valid};
\draw[flow] (g0.east) -- (r1.west);
\draw[flow] (g0.east) -- (r2.west);
\draw[flow] (g0.east) -- (r3.west);
\node[dec, right=0.4cm of r2] (d2) {rule refuses ($\rval = 1$), widens to $\Ir{1}$;\\ screen triggers, widens to $\Ir{1}$: \textbf{same}};
\node[dec, right=0.4cm of r3] (d3) {without $c_2$: rule certifies, reports PLAIN;\\ screen still triggers on $c_3$: \textbf{differ}};
\end{tikzpicture}
\end{center}

Where they differ, the screen widens without need and the rule is right to certify.
\end{exemplo}

Appendix~B fixed the policies before the run.

\begin{definicao}[The M1b policies]
$\mathrm{STOP}_r$, $r \in \{1, 2\}$: certify if $\Gz$ commits to a set and $\rval > r$, then report
PLAIN; otherwise report the band $\Ir{r}$. SCREEN: retract each claim, re-close, recompute the
committed result; if any differs from $\Gz$'s, report $\Ir{1}$, else PLAIN. SCREEN2: the same over
every retraction of one or two claims, widening to $\Ir{2}$. The \emph{committed result} is the set,
or the verdict class when there is no set. Policy comparisons are scored on informative queries
\emph{with} a committed set; without one, the two rules read literally differ by construction.
\end{definicao}

One direction always holds: whenever the rule refuses, the screen fires.

\begin{formulabox}[Argument: refusal implies trigger, at the same depth]
\emph{Premise P}: every set the committed-set procedure returns for a graph is valid in that graph.
If $\mathrm{STOP}_r$ refuses on a query with set $Z$, there is $S$ with $|S| \le r$ and $Z$ invalid in
$G_S = \mathrm{Meek}(\Chat, \Kset \setminus S)$. Either $G_S$ commits to some $Z_S$, valid by P and
therefore different from $Z$, or $G_S$ commits to no set and the verdict changes. In both cases the
depth-$r$ screen triggers. The converse fails: the screen also triggers when the set changes and $Z$
stays valid.
\end{formulabox}

This is not a theorem proved here, and the paper should not present it as one. P is a property of
the implementation (the closed-form optimal set, and one criterion both building the set and judging
it, which the back-door oracle of Chapter~\ref{ch:table4} showed can fail). What carries it is the
measurement: on all 2,212 policy rows every refusal at budget $r$ is matched by a screen triggering at
depth at most $r$, zero exceptions.

\textbf{The predictions}, written after M1 was read: Q1, rule and screen differ on at most 1\,\% of
committed queries at $b = 1$; Q2, at $b = 2$ on Panel~A the two-claim rule covers more than the
one-claim screen, with the interaction ceiling at most 0.31; Q3, at $b = 0$ the rule keeps the narrow
report; Q4, it refuses fewer than half of correct analyses; Q5, it covers at least 0.95 at $b = 1$.

\prompt[conceito=screen-leave-k-out,dificuldade=0.75]{State the premise on which ``if the stop rule
refuses, the screen triggers'' depends, and say why the converse fails.}{Premise: every set the
committed-set procedure returns is valid in the graph it was returned for. With it, if $Z$ is invalid
in some $G_S$ with $|S| \le r$, the committed result of $G_S$ is either a valid set (so different from
$Z$) or a no-set verdict, and the screen triggers. The converse fails because the screen also triggers
when a retraction changes the optimal set while $Z$ stays valid; there the rule certifies correctly
and the screen widens without need.}

\begin{antes}
\antesq[regra-de-parada]{With correct knowledge, on what share of informative committed analyses
does the stop rule refuse?}
\antesq[estrato]{What moves that number from 0.91 to 0.64, or to zero?}
\end{antes}

\section{The price of the guarantee: 0.910}

\begin{medido}[Refusal at $b = 0$, \texttt{results/interval/TABLE1\_L95.md}, M1b section]
Share of informative queries with a committed set on which the policy widens, knowledge correct,
network-bootstrap 95\,\% interval.

\medskip
\begin{tabular}{@{}lrrrr@{}}
panel (queries, networks) & $\mathrm{STOP}_1$ & $\mathrm{STOP}_2$ & SCREEN & SCREEN2 \\
\midrule
A (134, 20) & 0.910 [0.76, 1.00] & 0.955 & 0.910 & 0.963 \\
B (7, 3)    & 0.857 [0.80, 1.00] & 1.000 & 0.857 & 1.000 \\
\end{tabular}
\end{medido}

\textbf{Q4 failed, by a wide margin.} With every claim true, the stop rule refuses nine informative
committed analyses in ten. Where the knowledge carries identification, a single true claim, retracted,
breaks the set: the same table reports $\rval = 1$ on 0.910 of these 134 queries, which is the refusal
rate itself. The one-claim screen refuses on exactly the same queries. The left side of
\figref{fig:stop} shows the four rates.

\begin{figure}[!ht]
\centering
\includegraphics[width=\linewidth]{stop_refusal.pdf}
\caption{Left: with correct knowledge the stop rule and the one-claim screen refuse the same nine
analyses in ten; the 0.91 does not measure how demanding the rule is, it measures that the conclusion
rests on one claim. Right: at two wrong claims the two-claim rule covers more than the one-claim
screen and exactly as much as the pair screen.}
\label{fig:stop}
\end{figure}

The denominator must travel with the number: on committed queries outside the informative stratum
the rule refuses none (57 on Panel~A, 11 on Panel~B), and over all committed queries it refuses 0.639
(A) and 0.333 (B). \textbf{Q3 holds on the letter and is empty}: $\Lnf(\mathrm{STOP}_1)$ is below
$\Lnf(\Ir{1})$ in all four cells by 0.0003 to 0.0011, because read literally the rule widens on 97 to
99\,\% of informative queries. \textbf{Q5 holds}: at $b = 1$ it covers 0.979 and 0.979 on Panel~A
($n = 1000$, $20\,000$), 0.980 and 0.966 on Panel~B.

\prompt[conceito=regra-de-parada,dificuldade=0.7]{Give the three refusal rates of $\mathrm{STOP}_1$ at
$b = 0$ on Panel~A by denominator, and say which one Q4 was about.}{0.910 over the 134 informative
queries with a committed set (20 networks), the population of Q4, which predicted below 0.5; 0.639
over all committed queries; zero over the 57 committed queries outside the informative stratum, where
the class already identifies the effect.}

\begin{antes}
\antesq[screen-leave-k-out]{How much more does the two-claim rule cover than the one-claim screen at
$b = 2$, at which sample size, and on how many queries?}
\antesq[orcamento-retracao]{Where does a third of that difference come from?}
\antesq[screen-leave-k-out]{Against the pair screen, what is the difference?}
\end{antes}

\section{The starred comparison separates, then ties}

The comparison fixed on 7 September, $\mathrm{STOP}_2$ against SCREEN at $b = 2$ on Panel~A (59
committed informative queries, 8 networks), separates as Q2 predicted: $+0.256$ $[+0.149, +0.346]$ in
coverage at $n = 1000$ (0.986 against 0.731), $+0.269$ $[+0.163, +0.361]$ at $n = 20\,000$, and
$+0.305$ $[+0.182, +0.410]$ in the population (1.000 against 0.695). The interaction ceiling is 0.237
(A) and 0.286 (B). Q2 is supported as written, and two facts qualify it.

\begin{table}[!ht]
\centering\footnotesize
\caption{Where the $\mathrm{STOP}_2 - \mathrm{SCREEN}$ coverage gap comes from, Panel~A, $b = 2$,
$n = 1000$, committed informative queries. Source: \texttt{results/interval/TABLE1\_L95.md}, M1b
section.}
\label{tab:m1b-decomp}
\begin{tabular}{@{}lrrrr@{}}
\toprule
stratum & queries & share of replicates & gap inside & contribution \\
\midrule
interaction ($\rval = 2$, or a pair changes the result) & 14 & 0.24 & $+0.689$ & $+0.164$ \\
$\rval = 1$ (both policies refuse)                       & 34 & 0.58 & $+0.159$ & $+0.092$ \\
rest                                                     & 11 & 0.19 & $+0.005$ & $+0.001$ \\
\bottomrule
\end{tabular}
\end{table}

\textbf{About a third of the gap is the fallback, not the depth.} On the 34 queries with $\rval = 1$
both policies refuse, but the rule reports $\Ir{2}$ and the screen $\Ir{1}$, and with two wrong claims
$\Ir{2}$ covers more: $+0.092$ of the $+0.256$ (Table~\ref{tab:m1b-decomp}). Only the 14 interaction
queries measure depth. On Panel~B (14 queries, 4 networks) the gap is $+0.075$ $[+0.000, +0.175]$ at
$n = 1000$, all of it from 4 interaction queries, with an interval touching zero.

\textbf{A free screen at depth two ties exactly.} $\mathrm{STOP}_2$ against SCREEN2, both falling back
to $\Ir{2}$: coverage $+0.000$ $[+0.000, +0.000]$ at $n = 1000$, $-0.002$ $[-0.003, +0.000]$ at
$n = 20\,000$, zero on Panel~B. They decide differently on 5 of the 754 committed informative queries of
the five conditions ($141 + 162 + 73 + 161 + 217$), all of the example's class. The pre-fixed
comparison separated only because its free baseline stopped at depth one.

\textbf{Q1 failed by one query.} At $b = 1$, rule and screen differ on 2 of 162 committed queries
(1.23\,\%), both of the expected class: \texttt{Kampen\_2014}, ALN on PER, $Z = \{\mathrm{AFF},
\mathrm{SAN}\}$, $\rval = 3$; and \texttt{Schipf\_2010}, A on S, $Z$ empty, $\rval = 2$, where one
retraction makes the closed form return a non-empty optimal set for which the empty set is still valid.
Q2 and Q5 are supported, Q3 on the letter only, Q1 and Q4 falsified.

\prompt[conceito=orcamento-retracao,dificuldade=0.8]{Why does the $+0.256$ advantage of the two-claim
rule over the one-claim screen not measure only the value of certifying at depth two? Give the sample
size and the population line.}{Because the policies differ in two things: the depth of the certificate
and the budget of the fallback interval. On the 34 queries with $\rval = 1$ both refuse, but the rule
reports $\Ir{2}$ and the screen $\Ir{1}$; with two wrong claims $\Ir{2}$ covers more, and that stratum
contributes $+0.092$. Only the 14 interaction queries ($+0.164$) measure depth, and against a screen
that also tries pairs, with the same fallback, the difference is $+0.000$. The $+0.256$
$[+0.149, +0.346]$ is at $n = 1000$ on 59 queries and 8 networks; the population line of the same
contrast is $+0.305$ $[+0.182, +0.410]$.}

\begin{antes}
\antesq[frases-proibidas]{What does the pre-declared sentence turn the paper into?}
\antesq[regra-de-parada]{What survives M1b, measured?}
\antesq[screen-leave-k-out]{Which column is missing from the $\tau_b$ table of your draft?}
\end{antes}

\section{What survives, and the column your ranking table lacks}

Said without softening: \textbf{the sentence pre-declared on 7 September applies}. The starred
comparison separated against the one-claim screen and tied against the pair screen; at no depth we
measured does the stop rule compute something a leave-$k$-out screen does not. Nothing in the paper
may suggest that the radius, the certificate or the stop rule beats such a screen.

\begin{sonorio}[The tie, in the licensed form]
This is a measurement of when background-knowledge error is silent and what it costs. At every depth
we measured, the validity half of the sensitivity report matches a leave-$k$-out re-closure screen,
and the two differ only where a retraction changes the adjustment set while leaving it valid. We state
this as a property of the measure; what the report adds is the depth that was needed and the named
claims that witness it.
\end{sonorio}

What survives is measured, each number with its population. One wrong claim costs the committed
report 37 points of coverage (Panel~A, $n = 1000$: 0.967 to 0.597; 42 in the population), and the
knowledge interval restores coverage at 76 to 97\,\% of the class width ($\Lnf(\Ir{1})$ at $b = 1$,
0.758 to 0.974 over panels and sample sizes). Where an adjustment depends on knowledge, it almost
always rests on one claim (refusal 0.910; $\rnul$ at 1 or beyond 3 on 92.8\,\% of Panel~A's
established effects), so the report's most useful output is the \emph{name} of that claim; on the
ledger (a different harness: primary supplier, 157 exact rows, oracle CPDAG) the named claim was false on
58 of the 62 queries whose set was invalid at the truth (\texttt{results/ledger/CHECKS.md}, Table~8 in
Chapter~\ref{ch:frame}). The outline of 13 September
lists this framing as decision D9 and recommends accepting it rather than searching for a population
where conclusions rest on more claims, which would be selected for the answer. I recommend the same.

\textbf{The baseline your $\tau_b$ table is missing.} The ranking table of your draft
(\texttt{tab:tau} in \texttt{sections/06\_simulation.tex}) sets $\tau_b(\rval)$ against $\tau_b(\mathrm{SHD})$ and
$\tau_b(|\Kset|)$, two weak baselines. A reviewer who knows the area proposes the screen: the depth at
which leave-$k$-out re-closure first changes the set. By the argument above it never exceeds $\rval$
and equals it except where the set changes while staying valid (5 of 754 queries in M1b). The column
is unrun; your corruption states are stored, so it is cheap. I expect $\tau_b(\text{screen depth})$
close to $\tau_b(\rval)$; that is a prediction, not a measurement.

\begin{noprato}[Before the ranking table goes into the merged paper]
Add $\tau_b(\text{screen depth})$ on the same stored states and bootstrap. If it ties, state the
ranking as a property shared with the screen and say what the radius adds: the depth, the witness, the
certificate that no failure lies nearer.
\end{noprato}

\naopasse[orcamento-retracao]{Rebuild the decomposition of the starred comparison: queries, contribution
and mechanism per stratum. If you cannot say why the 34 queries with $\rval = 1$ contribute although both
policies refuse there, reread the paragraph on the fallback.}

\subsection*{Exercises}

\begin{exercicio}[metodo=estimation,conceito=regra-de-parada]
Before looking at the refusal table, you are given only the identity table for correct knowledge
($b = 0$), Panel~A, 134 informative committed queries: $\rval = 1$ on 0.910 of them and $\rval = 2$
on 0.045; no query differs between $\mathrm{STOP}_1$ and SCREEN. (a) Estimate the refusal rates of
$\mathrm{STOP}_1$, $\mathrm{STOP}_2$ and SCREEN. (b) Bound SCREEN2's from below, and say what the gap
between SCREEN2 and $\mathrm{STOP}_2$ can consist of. (c) Describe the population on which Q4
(refusal below 0.5) could have held.
\end{exercicio}
\begin{solucao}
(a) $\mathrm{STOP}_1$ refuses iff $\rval \le 1$: 0.910 (measured 0.910); $\mathrm{STOP}_2$ iff
$\rval \le 2$: $0.910 + 0.045 = 0.955$ (measured 0.955); SCREEN equals $\mathrm{STOP}_1$, no query
differing: 0.910. (b) Refusal implies trigger at the same depth, so SCREEN2 refuses at
least wherever $\mathrm{STOP}_2$ does: $\ge 0.955$. The excess can only be queries where a pair
retraction changes the set while $Z$ stays valid to depth two. Measured 0.963, so
$(0.963 - 0.955) \times 134 \approx 1$ query. (c) One where committed sets mostly survive every single
retraction, conclusions resting on two or more claims (larger chain components, denser knowledge near
the query). Here that is about one query in twenty; choosing such a population now would select for the
answer.
\end{solucao}

\begin{exercicio}[metodo=diagnosis,conceito=screen-leave-k-out]
A draft paragraph reads: \emph{``Our radius-based stop rule, which certifies the adjustment set up to
two retractions, improves coverage over leave-one-out screening by 25.6 points (0.986 vs.\ 0.731),
showing that the breakdown radius captures interactions between knowledge errors that a screen
misses.''} Diagnose every defect, with the number that exposes it, and say which sentence of
Chapter~\ref{ch:gamma} or this chapter should replace it.
\end{exercicio}
\begin{solucao}
(1) No population or sample size: $+0.256$ $[+0.149, +0.346]$ is Panel~A, $n = 1000$, 59 committed
informative queries on 8 networks; the population line is $+0.305$, and Panel~B's $+0.075$ touches zero.
(2) Depth is confounded with fallback: $+0.092$ comes from the 34 queries with $\rval = 1$ where both
policies refuse and differ only in widening to $\Ir{2}$ or $\Ir{1}$. (3) The conclusion is false at
matched depth: against the pair screen the difference is $+0.000$ $[+0.000, +0.000]$, with 5 of 754
queries decided differently, none unsafe. (4) ``Improves over'' is the forbidden claim that the radius
beats a leave-$k$-out screen. Replace the paragraph with this chapter's purple box, keeping the
one-claim contrast only with its sample size and decomposition.
\end{solucao}

\section*{Chapter map}
\begin{center}
\begin{tikzpicture}[node distance=0.6cm and 1.1cm]
\node[cmapc] (des) {7 Sep design:\\ policy rows};
\node[cmap, right=of des] (stop) {stop rule\\ $\mathrm{STOP}_r$};
\node[cmap, right=of stop] (scr) {leave-$k$-out screen\\ SCREEN, SCREEN2};
\node[cmapc, right=of scr] (frase) {pre-declared sentence\\ applies: a measurement};
\node[cmapr, below=of des] (ref) {refusal 0.910\\ with correct claims};
\node[cmapg, below=of stop] (sep) {separates: $+0.256$\\ ($n = 1000$) vs SCREEN};
\node[cmapr, below=of scr] (tie) {ties: $+0.000$\\ vs SCREEN2};
\node[cmap, below=of frase] (tau) {owed: screen-depth\\ column in $\tau_b$ table};
\draw[cedge] (des) -- node[above]{asks for} (stop);
\draw[cedge] (stop) -- node[above]{same at depth 1} (scr);
\draw[cedge] (scr) -- node[above]{tie triggers} (frase);
\draw[cedge] (des) -- node[left]{price cell} (ref);
\draw[cedge] (stop) -- node[left]{at $b = 2$} (sep);
\draw[cedge] (sep) -- node[above]{a third is fallback} (tie);
\draw[cedge] (frase) -- node[right]{implies} (tau);
\end{tikzpicture}
\end{center}
